\documentclass[10pt]{article}
\usepackage[margin=20mm]{geometry}
\usepackage{amsmath,amssymb,mathtools,fontspec,unicode-math,xurl,hyperref}
\IfFontExistsTF{DejaVu Serif}{\setmainfont{DejaVu Serif}}{\setmainfont{Latin Modern Roman}}
\setmathfont{latinmodern-math.otf}
\setlength{\parindent}{0pt}\setlength{\parskip}{6pt}\sloppy
\title{A proof-interface and normalization audit of
Strong Characterization for the Airy Line Ensemble}
\author{}
\date{}
\begin{document}
\maketitle
Abstract: Two deterministic tests show that the curvature/profile conditions in the central Section 3 reduction are impossible as printed. They are grouped as one Major proof-interface finding. The local curvature proof elsewhere uses the consistent constant, providing specific evidence for a repair; no counterexample to the final characterization is established. A separate Minor finding concerns an extra scale factor in the KPZ application. Probability quantifiers, affine mixtures, scaling, and selected conditional-limit and error-budget steps were also checked. All conclusions are restricted to the supplied arXiv version.\par
\section*{Identity, theorem scope, and result}
This report concerns Amol Aggarwal and Jiaoyang Huang, Strong Characterization for the Airy Line Ensemble [AH], assignment M25-0116. The inspected document is arXiv:2308.11908v2, 10 October 2025, 263 PDF pages. Page references below are PDF pages; the numbered text from page 3 uses the same numbers. The assignment’s journal DOI is 10.1007/s00222-025-01381-6 (https://doi.org/10.1007/s00222-025-01381-6). The journal PDF was not obtained, and equivalence to that version has not been established. A targeted primary-source check found the same v2 record and no separate correction notice; this is not an exhaustive absence claim.\par
Set \(a=2^{-1/2}\). Assumption 2.8 and Theorem 2.9 (p. 30) concern a continuous, \(\mathbb Z_{\ge1}\times\mathbb R\)-indexed Brownian Gibbs ensemble \(L\) satisfying \[\forall\varepsilon>0\ \exists K(\varepsilon)<\infty\
 \forall t\in\mathbb R:\quad
 \mathbb P\{|L_1(t)+at^2|\le\varepsilon t^2+K(\varepsilon)\}
 \ge1-\varepsilon.\] The probability bound is pointwise in \(t\), uniformly with the same constant; it is not a single simultaneous event for all \(t\). The asserted conclusion is \(L_j(t)=S_j(t)+lt+c\), with \((l,c)\) independent of the scaled parabolic Airy ensemble \[S_j(t)=a(A_j(t)-t^2).\] Stationarity and extremality are additional hypotheses for the deterministic-shift specialization.\par
The outcome is one Major and one Minor finding. The Major rating records a failed central intermediate claim and its actual use in the proof, not a demonstrated false main theorem. The inspected local proof strongly suggests coordinated corrections. The scope is a dependency audit with explicit analytic tests, not a complete independent reconstruction of every estimate in 263 pages. Independent scientific acceptance remains pending.\par
\section*{Problem 1: incompatible core curvature and profile conditions}
\subsection*{A deterministic contradiction in Lemma 3.21}
Equation reference: eq:incompatible. On p. 38, Lemma 3.21(3), equation (3.18), requires, for fixed \(\delta\in(0,1/4)\) and sufficiently large \(n\), a twice-differentiable \(h:[-T,T]\to\mathbb R\), where \[T=\delta^{-1}n^{1/3},\qquad b=2^{-7/6}(3\pi)^{2/3},\] satisfying, among its requirements, \[
 \sup_{|t|\le T}|h''(t)+a|\le\delta,\qquad
 \sup_{|t|\le T}|h(t)+at^2+bn^{2/3}|
       \le\frac{\delta n^{2/3}}2.\] The original page image confirms the negative exponent in the curvature constant. The same constant appears in the related Theorem 3.16, Assumption 3.18, and the application on p. 42.\par
Take the allowed value \(\delta=1/8\), and any integer \(n\ge\max\{2,N_0(\delta)\}\). Put \(g(t)=h(t)+at^2+bn^{2/3}\). The first bound in [eq:incompatible] (\#eq:incompatible) implies \(g''\ge a-\delta>0\). Thus \(g(t)-(a-\delta)t^2/2\) is convex, and its midpoint inequality gives \[g(T)+g(-T)-2g(0)\ge(a-\delta)T^2.\] The second bound instead gives \[g(T)+g(-T)-2g(0)\le2\delta n^{2/3}.\] Consequently \[a\le\delta+2\delta^3=\frac{33}{256},\] impossible since \(a=1/\sqrt2>1/2\). The argument only uses twice differentiability, exactly as stated; continuity of \(h''\) has not been added. It is deterministic and does not invoke an unproved stochastic rigidity estimate.\par
The underlying ensemble hypotheses are nonempty. Take the paper’s standard \(L=S\). Lemma 2.5 (p. 29) supplies its unit-variance Brownian Gibbs property, and Lemma 2.6 (p. 30) makes \(S_1(t)+at^2\) stationary. A finite quantile of \(|S_1(0)|\) supplies \(K(\varepsilon)\) uniformly for each fixed \(t\). Thus this is an admissible instance of Assumption 2.8. Nevertheless no \(h\) can satisfy the two requirements, so the event asserted to have probability at least \(1-\delta=7/8\) is empty. The additional approximation requirement involving \(L_{n+1}\) cannot rescue that event.\par
\subsection*{The Airy-gap hypothesis has an empty profile event}
Equation reference: eq:profile. On p. 37, equation (3.16) in Assumption 3.18 prints \[
 \mathcal F_n(t)=
 \left\{\max_{1\le j\le n}
 \left|x_j(t)-b\bigl(n^{2/3}-j^2\bigr)-f_n(t)\right|
 \le\delta_n n^{2/3}\right\}.\] Both the ordinary image and a higher-resolution inspection confirm \(j^2\), rather than \(j^{2/3}\). Under the nonintersecting bridge law above the finite lower boundary \(f_n\), one has \(x_n(0)>f_n(0)\) almost surely. The \(j=n\) term therefore exceeds \(b(n^2-n^{2/3})\). For sufficiently large \(n\), this is greater than \(\delta_n n^{2/3}\), because \(0<\delta_n<1/4\). Hence \(\mathcal F_n(0)\) is empty. Equation (3.17), requiring its probability to be at least \(1-\delta_n\), cannot hold.\par
This is a failure of satisfiability of the stated hypothesis and of its claimed verification. It is not a counterexample to a theorem under satisfiable hypotheses. The profile consistent with the global law and Lemma 3.21 is \(b(n^{2/3}-j^{2/3})\). At \(j=n\) that corrected profile vanishes and creates no conflict with the lower boundary.\par
\subsection*{Core dependency, counterevidence, and repair limits}
The proof on pp. 38–42 uses curvature approximation and Lemma 3.23 to obtain Lemma 3.21, then verifies Assumption 3.18 and applies the Airy-gap theorem (Theorem 3.19). Together with Proposition 3.20, these are inputs to Theorem 2.9. The incompatible bounds occur in the main proof reduction. This supports the Major classification of one combined interface defect, rather than two independent issue counts or a list of harmless isolated symbol errors.\par
There is substantial counterevidence against interpreting the finding as a false final characterization. Proposition 17.2, equation (17.2) on p. 189, bounds \(h''+2^{1/2}\). The gluing estimate on p. 191 retains this constant. Equation (17.36) on p. 199 differentiates the quadratic profile and obtains curvature \(-2^{1/2}\) exactly. The comparison calculation on p. 209 and the gap-scale limit in Lemma 19.4 on pp. 210–211 agree with that normalization.\par
Replacing \(+2^{-1/2}\) by \(+2^{1/2}\) in the second-derivative conditions removes the first contradiction: the exact quadratic \(h_0(t)=-at^2-bn^{2/3}\) has the corrected curvature and zero profile error. Replacing \(j^2\) by \(j^{2/3}\) removes the second contradiction. Intermediate gap factors must also agree with \[\sqrt2(S_j-S_{j+1})=A_j-A_{j+1};\] some displays on pp. 38–39 use the reciprocal factor, whereas the stated gap theorem and the final conclusion use \(\sqrt2\).\par
These are verified local consistency repairs. They do not show that \(h_0\) approximates the required random curve, prove all probability estimates, or certify every corrected application in the manuscript. The demonstrated defect is the printed central derivation; a counterexample to Theorem 2.9 has not been obtained. Author adoption and a correction of the journal version are not recorded.\par
\section*{Problem 2: the KPZ corollary states a different limiting law}
Corollary 25.1 (p. 255) states convergence to \(2^{-1/2}S^{(q)}\), with \(q=2^{-1/6}\), whereas its proof concludes convergence to \(S^{(q)}\). Definition (2.3) gives \[S_j^{(q)}(t)=q^{-1}S_j(q^2t)
       =2^{-1/3}A_j(2^{-1/3}t)-\frac{t^2}{2}.\] Thus \(S_1^{(q)}(0)=2^{-1/3}X\), where \(X=A_1(0)\) has the nondegenerate Tracy–Widom GUE law.\par
At \(t=0\), the paper’s own equation (25.1) gives the limiting distribution function \(F_{\rm TW}(2^{1/3}z)\), exactly that of \(S_1^{(q)}(0)\). The extra factor in the corollary instead gives \(2^{-5/6}X\) and distribution function \(F_{\rm TW}(2^{5/6}z)\). These distributions differ: a finite, nonzero, nondegenerate real random variable cannot have the same law as a nontrivial positive contraction of itself, since iteration would concentrate its law at zero. The discrepancy changes a stated application.\par
The local correction is to remove the prefactor \(2^{-1/2}\). The centering signs in (25.1) and (25.2), on pp. 255–256, also need to be consistent with the negative parabola: the centered variables use \(+t^2/2\) and \(+2^{-1/2}q^3t^2\), respectively. The \(t=0\) contradiction above is independent of these sign slips. This is a Minor finding in a secondary application; the main characterization does not depend on it. The correction agrees with the proof’s stated limit. The model-dependent log-gamma scale has not been independently rederived from its external references.\par
\section*{Positive checks and bounded proof qualifications}
\subsection*{No hidden moment hypothesis for affine shifts}
An arbitrary almost surely finite random pair \((l,c)\), independent of \(S\), is compatible with the main top-curve assumption. Choose quantiles \(M,N\) so that \[\mathbb P\{|l|>M\text{ or }|c|>M\}<\varepsilon/2,\qquad
 \mathbb P\{|S_1(0)|>N\}<\varepsilon/2.\] For each fixed \(t\), with probability at least \(1-\varepsilon\), \[|S_1(t)+at^2+lt+c|
 \le N+M+M|t|
 \le\varepsilon t^2+N+M+\frac{M^2}{4\varepsilon}.\] No moment of the shifts is used. Affine addition preserves the Brownian bridge specification, and independent mixtures preserve that same specification.\par
The transformation \(q^{-1}S(q^2t)\) retains Brownian variance one, since \(q^{-2}q^2=1\); its quadratic coefficient is \(q^3/\sqrt2\). This verifies \(q=2^{1/6}\sigma^{1/3}\) in Corollary 2.11. The sharper sublinear bound forces \(l=\ell\) almost surely: after division by \(t\), the stationary finite marginal and \(c\) vanish in probability. Stationarity similarly forces \(l=0\). Extremality removes a random intercept. The bounded function \(c\mapsto\mathbb E[\arctan(S_1(0)+c)]\) is strictly increasing, so a nontrivial split of the intercept law gives two distinct stationary Gibbs mixtures, contradicting extremality.\par
The step removing an independent intercept from \(X+c\stackrel d=X\) is valid without moments. The characteristic function of \(X\) is nonzero near zero, forcing that of \(c\) to equal one there and hence \(c=0\) almost surely.\par
\subsection*{Curvature, constants, and conditioning}
The curvature gluing has sufficient asymptotic room: with \(\varpi=n^{-13/15}\) and \(\eta=\exp(-\sqrt{\log n})\), \(n^{2/3}\varpi=n^{-1/5}\), and \(\varpi\eta^{-2}\) is smaller than every negative power of \(\log n\). Fixed multiples of \(\delta^{1/6}+(\log n)^{-1/3}\) can be absorbed in \(\delta^{1/8}+(\log n)^{-1/4}\) after the appropriate choices.\par
Two further local qualifications have explicit repairs. The application on p. 205 claims a full-domain \(C^{50}\) bound from the \(C^{45}\) bound in Definition 17.9. The general variant following Lemma 8.29 (p. 98), with \(m=44\) and \(\gamma=1/20000<1/(5m^2)\), instead needs precisely \(C^{45}\); its exponent \(2/44-1+1/20000\) beats the required \(-9/10\) budget, including the stated rescaling of particle number. On p. 42 the tolerance must be chosen after \(B=B'(\delta)+C_3(\delta)\), small enough for \(c(B)\); an unspecified positive \(c(B)\) need not exceed \(\delta^{20}\). An auxiliary tolerance below both \(c(B)/2\) and \((\delta/4)^8\) is compatible with the global-law theorem.\par
Conditioning on the exterior makes bridge data deterministic. The sequential formulation of Theorem 3.19 supplies uniformity over admissible configurations by contradiction and diagonal selection. Lemma 4.27 explicitly provides uniformity in initial data. Lemma 19.4 yields \(\sigma_{\nu;t}\to\sqrt2\); fixing a tolerance, sending \(n\) to infinity, and then decreasing that tolerance is a valid order.\par
\subsection*{Affine limits and independence}
In Section 20, \(N=n^{15}\), \(T=n^{1/3}\), and the endpoint normalization \((\xi_n,\zeta_n)\) is exterior-measurable. The deterministic sequential convergence of Proposition 20.1 likewise gives uniform conditional convergence for each bounded continuous test function on a fixed compact restriction.\par
The Cauchy argument on p. 216 needs a union of the two endpoint failure events, rather than their printed intersection. For fixed \(n\), convergence as the comparison index tends to infinity and the Airy concentration bound give a summable error \(C_\delta e^{-c_\delta n}\). Applying the union bound first over a finite tail of indices, and then taking limits, controls the supremum of all pairwise tail differences. This proves almost-sure Cauchy convergence of \((\xi_n,\zeta_n)\), rather than merely convergence in probability.\par
For bounded continuous \(F,G\), uniform conditional convergence gives \[\left|\mathbb E[F(L^n)G(\xi_n,\zeta_n)]
 -\mathbb E[F(S)]\,\mathbb E[G(\xi_n,\zeta_n)]\right|
 \le \|G\|_\infty\varepsilon_n+
 2\|F\|_\infty\|G\|_\infty\,\mathbb P(E_n^c),\] with both errors tending to zero. Dominated convergence then factorizes the limiting law. Taking \(S=L-\ell_\infty t-c_\infty\) on the same probability space gives the simultaneous pathwise decomposition and independence from the pair \((\ell_\infty,c_\infty)\). It does not assert that the two coefficients are independent of each other.\par
The coverage ledger records other locally repairable range and display mismatches in pp. 40 and 213–221, including the available stronger error bound in Lemma 20.2. They do not establish an independent obstruction to this mechanism. These bounded checks do not replace a full corrected proof.\par
\section*{Supporting checks}
Inspected dependency | Result and boundary of the check\par
Definitions, pp. 27–31, 43 | Pointwise probability, affine mixtures, scaling, and extremality checked.\par
Central reduction, pp. 34–42 | Explicit curvature contradiction and empty profile event; coordinated repairs identified.\par
Curvature and gap interfaces, pp. 95–98, 189–212 | Selected gluing, regularity, constants, conditioning, and limit order checked.\par
Affine limits, pp. 212–222 | Cauchy limits and independence audited; local corrections retained as qualifications.\par
Applications, pp. 255–256 | KPZ law checked; external log-gamma normalization not rederived.\par
Definition 2.2 reverses the upper/lower boundary labels relative to Definition 2.1. The surrounding ordering and Gibbs specification make the intended meaning clear; that isolated notation slip is not a separate finding. The proofs of all concentration, PDE, and monotonicity inputs in the intervening chapters and appendices have not been independently reconstructed. Unexamined steps are not classified as errors.\par
The exact supplied PDF SHA256 is\par
`c5a78895b4e686d1014fb3c2cc79a6eda800d41598f68c803866f16685ebf7e7`\par
Every report page is rendered for actual visual inspection. No author code, shell escape, new installation, or external model review was used. Independent scientific acceptance remains pending.\par
9 A. Aggarwal and J. Huang, Strong Characterization for the Airy Line Ensemble, arXiv:2308.11908v2, 10 October 2025, 263 pp. https://arxiv.org/abs/2308.11908v2. All numbered claims and equations cited here refer to the supplied PDF of this version.\par
\end{document}
